\documentclass[11pt,a4paper]{article}
\usepackage[margin=26mm]{geometry}
\usepackage{amsmath,amssymb,amsthm,booktabs,graphicx}
\usepackage{xcolor}
\usepackage[colorlinks=true,linkcolor=blue!45!black,citecolor=blue!45!black,urlcolor=blue!45!black]{hyperref}
\usepackage{microtype}
\newtheorem{proposition}{Proposition}
\newtheorem{theorem}{Theorem}
\newcommand{\E}{\mathbb{E}}
\newcommand{\dd}{\mathrm{d}}
\setlength{\emergencystretch}{2em}
\title{A finite noise window for concentration of a shared flow\\[0.4em]
\large Deterministic disturbance and steps toward selection}
\author{David Galbraith}
\date{Revised draft, 10 September 2026}
\begin{document}
\maketitle
\begin{abstract}
Can unequal sharing of a driven flow contribute to a physical account of
selection? For a bounded reflected diffusion in which channels share a conserved
total, we derive the exact joint stationary law and distinguish the population
of channel states from the flow they carry. With flow weight $g^m$ and diffusion
factor $g^q$, the population fraction above the pooled half-flow threshold tends
to zero as both state range and population grow precisely when $1\le q\le m+1$,
for exact powers. At fixed population $N$ it instead approaches $1/(2N)$ inside
the window. The upper edge is the established weighted-moment threshold, made
exact here for the coupled model. Independent simulations check the stationary
predictions on bounded grids. We distinguish this sorting result from spatial
channel formation, persistent structures, reproduction and portable heredity.
Existing and constructive models identify possible steps between them; they do
not yet establish the whole sequence in one physical system. A conditional
branching threshold states what further measurements would be needed to predict
lineage establishment. Empirical applicability remains untested.
\end{abstract}

\section{The physical question and the upper boundary}

The motivating proposal is conditional: disturbance associated with a driven
flow could help produce concentrated sharing only within a suitable scaling
window, and that concentration could become relevant to selection if states
persist and are copied. The window is essential to this proposal. A mechanism
that concentrates under every disturbance rule would not establish it.

The research programme asks whether this sequence can be constructed without
adding random forcing as a fundamental ingredient: deterministic fine-scale
dynamics produce effective variation; variation affects the division of flow;
and inherited states use some of that flow to reproduce. This paper separates
an exact stationary result from computational thought experiments that explore
those connections. The proposed synthesis is not equivalent to the known
power-law calculation on which one step depends.

A collection of paths can be counted in two ways: by the number of paths in each
state, or by the fraction of throughput they carry. A rare large state can matter
little to the first count and greatly to the second. The question here is when
that difference becomes extreme as the accessible state range grows.

Noise affects this comparison through residence times. In a scalar driftless
It\^o diffusion with reflecting boundaries, the stationary density is inversely
proportional to diffusivity. State-dependent noise and occupation probabilities
have a long history \cite{landauer,maes}. Weighting a power-law distribution by a
power of its argument introduces a standard moment threshold \cite{newman}.
Neither ingredient is new.

Newman explicitly discusses both the general moment condition and the fraction
of a weighted total held in a population tail (Secs.~III.B and III.D of
Ref.~\cite{newman}). Their use for a fixed shared flow requires care: each channel's share depends on
every other channel, and the same normalization can change the physical clock
of the state dynamics. We give an exactly solvable coupled model in which this
dependence is explicit. The result identifies the finite-population correction
and makes the observable behind the noise window precise.

For $m=1$, the upper edge is $q=2$. Beyond it, increasing the upper state cutoff
does not make a vanishing population fraction account for half the pooled flow.
This does not mean that flow becomes equally distributed. A small but nonzero
fraction can still carry half. Nor is $q$ a noise amplitude: it describes how
diffusivity changes with state. Multiplying every diffusivity by one constant
changes relaxation times, but leaves all stationary results below unchanged.

\section{A model with conservation imposed at every instant}

Let $N\geq1$ channels have dimensionless states $g_i\in[1,R]$, where $R>1$.
The model specifies possible flow routes and their response law; it does not
prescribe which routes become dominant. The lattice experiment in
Sec.~\ref{sec:thought} tests differentiation from identical coarse-route states.
Spatial formation and connectivity of routes in a material are abstracted out:
no spatial energy-gradient field or changing network topology is modeled here.
For $m>0$, define the total weight $S$ and the actual flow through channel $i$ by
\begin{equation}
 S(\mathbf g)=\sum_{j=1}^N g_j^m,\qquad
 J_i(\mathbf g)=I\frac{g_i^m}{S(\mathbf g)}.
 \label{eq:share}
\end{equation}
Thus $\sum_i J_i=I$ for every configuration, not only in expectation. We choose
the interior dynamics
\begin{equation}
 \dd g_i=\sqrt{2D_i(\mathbf g)}\,\dd W_i,\qquad
 D_i(\mathbf g)=\frac{N g_i^q}{S(\mathbf g)},
 \label{eq:model}
\end{equation}
with independent Wiener processes, normal reflection at both endpoints, and
zero interior drift. Boundary local-time terms are implicit. A positive common
prefactor has been absorbed into the time unit. The lower cutoff sets the state
unit; all powers refer to this specified coordinate.

The normalization in Eq.~\eqref{eq:model} is a constitutive assumption. On the
line $q=m$, $D_i=(N/I)J_i$, so noise variance rate is proportional to actual
flow. Away from that line the model keeps the common factor $1/S$. The apparently
similar choice $D_i\propto J_i^{q/m}$ instead contains $S^{-q/m}$ and defines
different dynamics. Conservation of flow alone does not select either rule.

The It\^o convention would follow from an appropriate small-kick limit of
conditionally unbiased additive events in $g$, evaluated at the pre-event
state. Such a limit also needs control of temporal correlations and higher
moments. Unresolved deterministic dynamics do not automatically satisfy these
conditions. Under a nonlinear change $y=h(g)$, It\^o's formula adds the drift
$D(g)h''(g)$ even when the drift in $g$ is zero. The zero-drift assumption must
therefore be justified in the measured coordinate, rather than restored after
arbitrary relabelling.

\section{Exact stationary population and flow distributions}

On $[1,R]$, write
\begin{equation}
 Z_R=\int_1^R g^{-q}\dd g,\quad p_0(g)=\frac{g^{-q}}{Z_R},\quad
 \mu_R=\int_1^R g^m p_0(g)\dd g,\quad
 f_R(g)=\frac{g^m p_0(g)}{\mu_R}.
 \label{eq:base}
\end{equation}
All quantities are finite at finite $R$, for every real $q$.

\begin{proposition}
The unique stationary density of Eq.~\eqref{eq:model} on the reflecting box is
\begin{equation}
 \pi_{N,R}(\mathbf g)=\frac{S(\mathbf g)}{N\mu_R}
                 \prod_{j=1}^N p_0(g_j).
 \label{eq:joint}
\end{equation}
Its tagged-channel marginal and its pooled flow density are, respectively,
\begin{align}
 p_{N,R}(g)&=\left(1-\frac1N\right)p_0(g)+\frac1N f_R(g),
 \label{eq:marginal}\\
 \E_{\pi}\!\left[\sum_{i=1}^N\frac{J_i}{I}\delta(g-g_i)\right]&=f_R(g).
 \label{eq:flow}
\end{align}
\end{proposition}

\begin{proof}
Integration of $S\prod_j p_0(g_j)$ gives $N\mu_R$, establishing normalization.
For each coordinate,
\[
 D_i\pi_{N,R}=\frac{1}{Z_R\mu_R}\prod_{j\ne i}p_0(g_j)
\]
is independent of $g_i$. Every probability current
$-\partial_i(D_i\pi_{N,R})$ therefore vanishes, including at the boundaries.
The diffusion is uniformly nondegenerate on the connected compact box, giving
uniqueness. Expanding the sum $S$ and integrating out the other coordinates
gives Eq.~\eqref{eq:marginal}. Multiplication by $J_i/I=g_i^m/S$ cancels $S$
in the joint law; integration and summation give Eq.~\eqref{eq:flow}.
\end{proof}

The calculation also has a simple clock interpretation. If
$L_0=\sum_i g_i^q\partial_i^2$ is the generator of independent reflected
channels, the coupled generator is $L=(N/S)L_0$. Configurations with larger
$S$ are traversed more slowly relative to that independent process; their
stationary weight acquires the factor $S$. This explains both the solvability
and the dependence of the result on the chosen normalization.

An exact independent sampling construction is useful: choose one channel
uniformly, draw its state from $f_R$, and draw the other $N-1$ states independently
from $p_0$. Averaging over the chosen channel reproduces Eq.~\eqref{eq:joint}.
That designation is resampled for each configuration; it does not identify a
permanent winner.

The correction matters even in simple cases. For $N=1$ and $q=m$, the channel
carries all the flow and has constant diffusivity. Its stationary density is
uniform, not $g^{-m}$. At fixed $R$, the tagged marginal approaches $p_0$ as
$N\to\infty$. The independent-channel calculation is recovered in that limit.

\section{The concentration observable and its window}

Let $a_R$ be the state threshold above which half of the pooled flow lies:
\begin{equation}
 \int_{a_R}^R f_R(g)\dd g=\frac12.
\end{equation}
Define $P_0(R)=\int_{a_R}^R p_0(g)\dd g$ and
$P_N(R)=\int_{a_R}^R p_{N,R}(g)\dd g$. Equation~\eqref{eq:marginal} gives
\begin{equation}
 \boxed{P_N(R)=\left(1-\frac1N\right)P_0(R)+\frac{1}{2N}.}
 \label{eq:correction}
\end{equation}
This counts a population fraction above one threshold fixed across stationary
observations. It is a pooled observable, with each observation weighted equally
in physical time.

\begin{theorem}
For Eqs.~\eqref{eq:share}--\eqref{eq:model} with exact powers and $m>0$,
$P_N(R)\to0$ as both $R\to\infty$ and $N\to\infty$ if and only if
\begin{equation}
 \boxed{1\leq q\leq m+1.}
 \label{eq:window}
\end{equation}
At fixed $N$, the limit inside this window is $1/(2N)$.
\end{theorem}

\begin{proof}
Put $A=m-q+1$ and $B=1-q$. Direct integration gives
\begin{equation}
 a_R=\begin{cases}
 [(R^A+1)/2]^{1/A},&A\ne0,\\
 \sqrt R,&A=0,
 \end{cases}\qquad
 P_0(R)=\begin{cases}
 (R^B-a_R^B)/(R^B-1),&B\ne0,\\
 \log(R/a_R)/\log R,&B=0.
 \end{cases}
 \label{eq:quantile}
\end{equation}
For $q<1$, $a_R/R\to2^{-1/A}$ and
$P_0\to1-2^{-B/A}>0$. At $q=1$,
$P_0\sim\log2/(m\log R)\to0$. For $1<q<m+1$, $a_R$ grows proportionally
to $R$, while the normalized population tail above it vanishes.
At $q=m+1$, $P_0=1/(R^{m/2}+1)\to0$.
For $q>m+1$, $a_R\to2^{1/(q-m-1)}$, so
$P_0\to2^{-(q-1)/(q-m-1)}>0$.
Equation~\eqref{eq:correction} supplies the finite-$N$ result and the joint limit.
\end{proof}

Either iterated limit gives this classification, as does any sequence with both
$R$ and $N$ tending to infinity. Stationarity is taken at each bounded pair
$(N,R)$ first. No unbounded diffusion or uniform equilibration time is asserted.
The same classification holds for any fixed pooled flow fraction strictly
between zero and one, with the corresponding fraction divided by $N$ replacing
$1/(2N)$.

The integrals behind the two edges are $\int g^{-q}\dd g$ and
$\int g^{m-q}\dd g$. At the lower edge the population normalization is
logarithmic; at the upper edge the flow normalization is logarithmic.
The upper-edge population fraction itself decays as a power of $R$.
Figure~\ref{fig:window} displays the finite-range expressions.

\begin{figure}[t]
\centering
\includegraphics[width=0.94\textwidth]{window_figure.pdf}
\caption{Exact stationary predictions, not empirical measurements or large-range
dynamic simulations. The asymptotic window is between the dotted lines. Finite
range smooths both boundaries; finite population supplies a floor. $P_0$ is the
large-population limit at each fixed $R$. The curve at $R=10^8$ is not the
infinite-range limit.}
\label{fig:window}
\end{figure}

\subsection{Pooled concentration is not a snapshot winner count}

In a single configuration, rank shares in descending order and count the
smallest number of whole channels whose shares sum to at least one half.
Call that count divided by $N$ $H_N$. It always satisfies $H_N\geq1/N$.
In contrast, $P_N$ can approach $1/(2N)$ because its fixed threshold need not
carry half the flow in every snapshot. These observables must not be substituted
for one another. The numerical checks measure $\E H_N$ separately against the
exact joint sampler. Neither statistic measures how long a particular channel
remains dominant. Stationary concentration alone establishes neither persistent
structure nor Darwinian selection.

For example, at $m=1$, $q=2$, $R=10^8$, $P_0\simeq0.000100$, whereas
$P_{40}\simeq0.012598$. This is a substantial relative correction to the tail
fraction, even though $1/N$ is small.

\subsection{Limiting slopes do not determine the endpoints}

The inclusive endpoints require exact powers or additional tail assumptions.
On $g\geq e$, consider the unnormalized flow weight $w(g)=g$ and noise factor
$d(g)=g^2(\log g)^2$. Their limiting log-log slopes are $m=1$ and $q=2$, but
$\int_e^\infty w(g)/d(g)\dd g$ converges. The upper-edge separation fails.
This example also fits the coupled construction with $S=\sum_i w(g_i)$ and
$D_i=N d(g_i)/S$: replace $p_0$ by the normalized $1/d$ and $f_R$ by the
normalized $w/d$ in the proof. Logarithmic corrections do not change the strict
interior power inequalities, but can decide the boundary cases. An estimated
slope near an edge is consequently insufficient to classify a measured system.

\section{What would make the result physically useful?}

The upper bound excludes a proposed mechanism when the measured disturbance
scaling is too steep relative to the flow weight, provided the rest of the model
applies. It also distinguishes state-dependent disturbance from monotonic
reinforcement intuition. At $m=1$, the large-range, large-population half-flow
fraction vanishes for $q=2$, but approaches $1/8$ for $q=2.5$.
Above the boundary it can still be very small, especially close to $q=2$.
The mathematical transition is the loss of a vanishing fraction, not an abrupt
return to equal sharing in a finite apparatus.

The basic inequality is a standard power-law moment condition \cite{newman}.
Its value here is the connection to actual conserved shares and the observable
correction in Eq.~\eqref{eq:correction}. These make a conditional prediction
that can be checked. They do not establish a new general law of nature, and no
priority claim for the elementary threshold is made.

One possible microscopic route to state-dependent variance uses contact patches.
If each patch contributes an unbiased additive perturbation in the same measured
coordinate, $n$ independent patches give variance proportional to $n$; perfectly
coherent patches within a channel give variance proportional to $n^2$.
Inferring $q=m$ or $q=2m$ additionally requires $n\propto g^m$, fixed coupling
per patch, and a suitable time scale. It does not derive the common $1/S$ clock
in Eq.~\eqref{eq:model}. In particular, a rule proportional to squared actual
flow would have a different normalization. Coherence within a channel must also
be distinguished from correlations between channels, which violate the
independent-noise assumption.

A fair physical test needs simultaneous channel states and individual flows,
with a fixed total and a stable channel population. It must estimate both
conditional drift and variance rate in the chosen coordinate, check cross-channel
correlations, and allow enough relaxation before comparing stationary statistics.
For Eq.~\eqref{eq:model}, conditioning on the configuration or accounting for $S$
is necessary: a regression on $g_i$ alone need not recover $q$. Drift, measurement
noise, jumps, changing populations and limited state range can each prevent the
test from addressing this model. Power-law fits themselves also need checks
against alternatives \cite{clauset}.

A useful experiment would vary the state dependence of disturbance across the
upper edge while measuring the finite-range prediction, rather than expecting
an infinite-limit discontinuity. No physical dataset is used as validation in
this paper.

\section{Thought experiments toward selection}
\label{sec:thought}

\subsection{Effective noise in a deterministic world}

The first construction, \texttt{sims/core/08}, puts each of forty coarse routes
over a block of a coupled logistic-map lattice. The lattice evolves
deterministically. A route responds to the difference between its block mean
and a slowly updated running mean, with a coupling proportional to its share.
Random numbers initialize the lattice and routes; none are drawn during their
subsequent evolution. Fixing the complete initial state fixes the trajectory.

In the retained default run, the chaotic lattice gives a final top-three flow
share of $0.62$, compared with $0.08$ for a frozen lattice. Its sampled block
means have mean lag-one autocorrelation about $0.01$, and the repeated run
reports identical summary statistics. These are demonstrations of deterministic
driving with irregular coarse fluctuations and concentrated flow, not a proof
that the coarse process is Markovian or converges to the It\^o model above.

The periodic control is informative: it gives an even larger top-three share,
$0.95$. The experiment therefore does not show that chaos is necessary, despite
an earlier description of that control. Low autocorrelation alone would not
establish unpredictability from all coarse information either. The contribution
of the lattice experiment is to make the absence of ongoing random input
explicit and reproducible within a flow-sharing example.

Deterministic dynamics admitting stochastic effective descriptions are already
established mathematically \cite{melbourne,gottwald}. Their role here is as a
candidate physical source for the disturbance, not a new discovery that
determinism and statistical noise can coexist. A diffusion limit for this
particular lattice, including its effective drift and stochastic convention,
remains to be derived. This matters because deterministic multiplicative
forcing need not produce a driftless It\^o limit \cite{gottwald}.

\subsection{Contact geometry and the window are separate tests}

The next construction, \texttt{sims/core/09}, recruits more lattice patches for
a route with a larger share, keeping coupling per patch fixed. In the retained
run, the variance-versus-mean-share slope is $1.01$ for sums over separate
patches. A common coherent patch gives slope $2.00$; using fixed patches with
amplitude proportional to share gives $1.99$. These are across-route diagnostic
slopes, not direct estimates of conditional infinitesimal diffusivity.

The final top-three shares are $0.17$, $0.08$ and $0.63$, respectively.
The coherent case uses the same fluctuation for all routes and so changes
cross-route correlation as well as variance scaling. It cannot isolate the
effect of moving $q$ to the upper edge. Both slopes near two are also at the
$m=1$ boundary, rather than beyond it. The lattice examples consequently
support a mechanism for effective disturbance and its scaling; they do not
demonstrate the full two-sided window in a single deterministic substrate.

The exact theorem answers a different, controlled question: for the specified
diffusion, when does pooled concentration vanish as a population fraction?
The lattice routes have a lower clipping rule and no fixed upper reflecting
cutoff, and their effective drift and shared clock have not been matched to
Eq.~\eqref{eq:model}. A next test must connect those models before treating
their results as one experiment.

\subsection{From unequal sharing to differential reproduction}

Concentration alone contains no lineage. Suppose each route also has a store
of usable work $A_i$, increasing at rate $\dot A_i=\eta J_i$, and a copy costs
$E_{\rm copy}$. With constant flow and no intervening loss, the time between
copies is $E_{\rm copy}/(\eta J_i)$. For time-varying flow, accumulated work
sets the number of completed copies, up to the remainder in the store. Thus
the ideal long-run rate is $b_i=\eta\overline J_i/E_{\rm copy}$.
This derives the rate conditional on a copying mechanism, an efficiency and a
cost; it does not derive the existence of copying machinery.

In \texttt{sims/core/18}, a funded copy replaces a randomly chosen route with
an inherited state plus a perturbation. The default run with copy cost $200$
ends at mean state $8.67\pm1.27$, compared with $4.32\pm1.17$ at cost $500$
and $1.05\pm0.19$ without copying. These are means and population standard
deviations across three runs, not confidence intervals. Fitted copying-rate
slopes are $0.0045$ and $0.0017$, against ideal values $0.0050$ and $0.0020$;
replacement can erase stored resources, and self-targeted attempts are counted.

Here the connection to selection becomes explicit: state varies, is inherited,
and affects reproduction through throughput. The simulation supplies
inheritance and replacement rules and uses random draws during evolution.
It is therefore a separate thought experiment, not a completed demonstration
of selection driven entirely by the deterministic lattice. Other repository
models investigate whether copying refreshes memory faster than restoring
dynamics erase it. Those address another necessary condition for the proposed
sequence, rather than following automatically from the concentration theorem.

\subsection{What the combined proposal commits to}

The scientific proposal is the connection between these steps, including the
bounded concentration condition. The theorem constrains this concentration
mechanism, not every possible route to Darwinian selection. It would gain substance from a construction
in which the same physical rules supply effective variation, cross both noise
boundaries, preserve inherited differences and fund reproduction. Whether
channels themselves must be supplied or can form from an initially unchannelled
material is a further question. No such complete construction is established
here. Nor does the prior art cited for the separate ingredients establish that
the complete proposed mechanism is already known.

All default outputs discussed in this section are retained with source hashes
and reproduction commands in \texttt{docs/research\_examples/}. They report
finite-time toy-model behaviour. They neither establish empirical validity nor
replace the stationary checks below.


\section{Numerical checks of the stationary result}

The repository implements a continuous-time nearest-neighbour chain on a uniform
state lattice with spacing $h$. Each allowed jump of channel $i$ by $+h$ or $-h$
has rate
\begin{equation}
 \lambda_i(\mathbf g)=\frac{N g_i^q}{S(\mathbf g)h^2}.
\end{equation}
Outward boundary jumps are absent. Exponential holding times and rate-weighted
event selection simulate this finite chain without a time-step approximation.
The lattice invariant law is exactly Eq.~\eqref{eq:joint} with sums replacing
integrals. Grid refinement remains necessary for a continuous-state comparison.
Observations are taken at fixed physical times, because sampling only jump
events would bias residence-time statistics.

Eight automated tests include detailed balance on every edge of an enumerated
three-channel, three-state system; independent marginal calculations; the
single-channel case; the finite-population quantile identity; sampling checks;
continuum integral comparisons; grid refinement; and relaxation from both
boundaries in a simple control. The retained ensemble runs use six exponent pairs:
$(1,0.5)$, $(1,1)$, $(1,1.5)$, $(1,2)$, $(1,2.5)$ and $(2,4)$.
For each pair they use lower-boundary, upper-boundary, log-uniform and exact
stationary initial conditions, with three seeds and 128 independent ensembles
per seed. Physical observation times are $0,1,4,16$.

\begin{table}[!ht]
\centering
\begin{tabular}{rrrrr}
\toprule
$N$ & Grid intervals & Occupancy CDF & Flow CDF & Mean snapshot fraction\\
\midrule
8 & 4 & 0.017014 & 0.021730 & 0.011423\\
8 & 8 & 0.020299 & 0.024502 & 0.005979\\
40 & 8 & 0.009496 & 0.008539 & 0.001585\\
\bottomrule
\end{tabular}
\caption{Largest final absolute errors across the six exponent pairs and four
initial conditions for each retained run, all on $[1,3]$. CDF errors compare
with the exact lattice law. The snapshot reference uses 100,000 independent
joint-law samples per pair. Seeds are pooled at the final observation time.}
\label{tab:checks}
\end{table}

All retained final errors are below the declared tolerance of $0.05$
(Table~\ref{tab:checks}). This is an acceptance criterion, not a confidence
interval. Separate exact-law calculations check grid refinement and evaluate
large state ranges; they are not dynamical runs at those ranges. The checks
support the implementation on the tested bounded grids. They do not establish
large-range equilibration times or empirical applicability.

The earlier independent-walker check is excluded from this evidence. It started
log-uniformly, which is already the target stationary law at $q=1$, and did not
establish relaxation. It is retained in the repository as a labelled legacy
diagnostic. The exact power integrals remain valid.

\section{Conclusion}

The conserved-flow model gives an exact concentration window and a
finite-population correction. The upper bound states when rare large states
cease to dominate pooled throughput; its mathematics is established. The
physical proposal is the connection from such sorting to structures whose
variations persist and reproduce. The repository makes several steps concrete
and keeps their assumptions separate. A complete account must derive the
connections under consistent physical rules and test them independently;
concentration alone is not a derivation of selection.

\paragraph{Code and reproducibility.}
Manuscript sources, model implementations, tests and retained data are at
\url{https://github.com/divadwg/selection-as-physics}.
Build and reproduction commands are in \texttt{paper/README.md}; the repository
separates current evidence from historical interpretations and secondary models.

\begin{thebibliography}{99}
\bibitem{landauer} R. Landauer, Motion out of noisy states,
\emph{J. Stat. Phys.} \textbf{53}, 233--248 (1988).
\href{https://doi.org/10.1007/BF01011555}{doi:10.1007/BF01011555}.
\bibitem{maes} C. Maes and K. Neto\v{c}n\'y, Heat bounds and the blowtorch theorem,
\emph{Ann. Henri Poincar\'e} \textbf{14}, 1193--1202 (2013).
\href{https://arxiv.org/abs/1207.1122}{arXiv:1207.1122}.
\bibitem{newman} M. E. J. Newman, Power laws, Pareto distributions and Zipf's law,
\emph{Contemporary Physics} \textbf{46}, 323--351 (2005).
\href{https://arxiv.org/abs/cond-mat/0412004}{arXiv:cond-mat/0412004}.
\bibitem{clauset} A. Clauset, C. R. Shalizi and M. E. J. Newman,
Power-law distributions in empirical data,
\emph{SIAM Review} \textbf{51}, 661--703 (2009).
\href{https://arxiv.org/abs/0706.1062}{arXiv:0706.1062}.
\bibitem{melbourne} I. Melbourne and A. M. Stuart,
A note on diffusion limits of chaotic skew-product flows,
\emph{Nonlinearity} \textbf{24}, 1361--1367 (2011), corrected version (2015).
\href{https://arxiv.org/abs/1101.3087}{arXiv:1101.3087}.
\bibitem{gottwald} G. A. Gottwald and I. Melbourne,
Homogenization for deterministic maps and multiplicative noise,
\emph{Proc. R. Soc. A} \textbf{469}, 20130201 (2013), corrected version (2015).
\href{https://arxiv.org/abs/1304.6222}{arXiv:1304.6222}.
\bibitem{branching} S. Hautphenne, G. Latouche and G. Nguyen,
Extinction probabilities of branching processes with countably infinitely many types,
\emph{Advances in Applied Probability} \textbf{45}, 1068--1082 (2013).
\href{https://arxiv.org/abs/1211.4129}{arXiv:1211.4129}.
\bibitem{pearson} J. E. Pearson, Complex patterns in a simple system,
\emph{Science} \textbf{261}, 189--192 (1993).
\href{https://arxiv.org/abs/patt-sol/9304003}{arXiv:patt-sol/9304003}.
\bibitem{spots} W. N. Reynolds, S. Ponce-Dawson and J. E. Pearson,
Self-replicating spots in reaction-diffusion systems,
\emph{Physical Review E} \textbf{56}, 185--198 (1997).
\href{https://doi.org/10.1103/PhysRevE.56.185}{doi:10.1103/PhysRevE.56.185}.
\bibitem{chemicalspots} K.-J. Lee, W. D. McCormick, J. E. Pearson and H. L. Swinney,
Experimental observation of self-replicating spots in a reaction--diffusion system,
\emph{Nature} \textbf{369}, 215--218 (1994).
\href{https://doi.org/10.1038/369215a0}{doi:10.1038/369215a0}.
\bibitem{peptide} D. H. Lee et al., A self-replicating peptide,
\emph{Nature} \textbf{382}, 525--528 (1996).
\href{https://doi.org/10.1038/382525a0}{doi:10.1038/382525a0}.
\bibitem{lbm} M. Junk, A. Klar and L.-S. Luo, Asymptotic analysis of the lattice Boltzmann equation,
\emph{Journal of Computational Physics} \textbf{210}, 676--704 (2005).
\href{https://doi.org/10.1016/j.jcp.2005.05.003}{doi:10.1016/j.jcp.2005.05.003}.
\end{thebibliography}
\end{document}
